
\documentclass{article}
\usepackage{colortbl}
\usepackage{makecell}
\usepackage{multirow}
\usepackage{supertabular}

\begin{document}

\newcounter{utterance}

\centering \large Interaction Transcript for game `cladder', experiment `full\_v1.5\_default', episode 4040 with refineonly{-}llama{-}3.1{-}8b{-}instruct.
\vspace{24pt}

{ \footnotesize  \setcounter{utterance}{1}
\setlength{\tabcolsep}{0pt}
\begin{supertabular}{c@{$\;$}|p{.15\linewidth}@{}p{.15\linewidth}p{.15\linewidth}p{.15\linewidth}p{.15\linewidth}p{.15\linewidth}}
    \# & \multicolumn{2}{c}{Player} && \multicolumn{2}{c}{Game Master} \\
    \hline

    \theutterance \stepcounter{utterance}  
    & & & \multicolumn{4}{p{0.6\linewidth}}{
        \cellcolor[rgb]{0.9,0.9,0.9}{
            \makecell[{{p{\linewidth}}}]{
                \texttt{\tiny{[P1$\langle$GM]}}
                \texttt{You are an expert in causal inference. The following question is not a typical commonsense query, but rather a meticulously designed question created by a professor specializing in causal inference, intended to assess the students' mastery of the course content.} \\
\texttt{Imagine a self{-}contained, hypothetical world with only the following conditions, and without any unmentioned factors or causal relationships: Smoking has a direct effect on effort and lactose intolerance. Effort has a direct effect on lactose intolerance. For nonsmokers, the probability of being lactose intolerant is 77\%. For smokers, the probability of being lactose intolerant is 37\%. Will smoking decrease the chance of being lactose intolerant?} \\
\texttt{Start your answer with 'Yes' or 'No', followed by additional reasoning or evidence to support your explanation.} \\
            }
        }
    }
     \\ \\

    \theutterance \stepcounter{utterance}  
    & & \multicolumn{4}{p{0.6\linewidth}}{
        \cellcolor[rgb]{0.95,0.95,0.95}{
            \makecell[{{p{\linewidth}}}]{
                \texttt{\tiny{[P1$\leftrightarrow$P]}}
                \texttt{{[}DRAFT{]}} \\
\texttt{No.} \\
\\ 
\texttt{To answer this question, we need to consider the causal relationships and their directions. Smoking has a direct effect on effort, and effort has a direct effect on lactose intolerance. } \\
\\ 
\texttt{Let's denote the variables as follows: } \\
\\ 
\texttt{{-} S: Smoking status (1 for smokers, 0 for nonsmokers)} \\
\texttt{{-} E: Effort (1 for high effort, 0 for low effort)} \\
\texttt{{-} L: Lactose intolerance (1 for lactose intolerant, 0 for lactose tolerant)} \\
\\ 
\texttt{Given the conditions, we can establish the following relationships:} \\
\\ 
\texttt{1. Smoking affects effort: P(E|S=1) > P(E|S=0)} \\
\texttt{2. Effort affects lactose intolerance: P(L|E=1) > P(L|E=0)} \\
\\ 
\texttt{We are also given the following probabilities:} \\
\\ 
\texttt{{-} P(L|S=0) = 0.77 (probability of being lactose intolerant for nonsmokers)} \\
\texttt{{-} P(L|S=1) = 0.37 (probability of being lactose intolerant for smokers)} \\
\\ 
\texttt{To determine the effect of smoking on lactose intolerance, we need to consider the indirect effect of smoking through effort. Since effort has a direct effect on lactose intolerance, we can use the following formula to calculate the total effect of smoking on lactose intolerance:} \\
\\ 
\texttt{P(L|S=1) = P(L|E=1) \textbackslash{}* P(E|S=1) + P(L|E=0) \textbackslash{}* P(E=0|S=1)} \\
\\ 
\texttt{We know that P(L|E=1) > P(L|E=0) and P(E|S=1) > P(E|S=0). However, we also know that P(L|S=1) = 0.37, which is less than P(L|S=0) = 0.77.} \\
\\ 
\texttt{To reconcile this, we need to consider the fact that the probability of being lactose intolerant for smokers (P(L|S=1)) is a weighted average of the probabilities of being lactose intolerant for individuals with high effort (P(L|E=1)) and low effort (P(L|E=0)), where the weights are the probabilities of having high effort (P(E=1|S=1)) and low effort (P(E=0|S=1)).} \\
\\ 
\texttt{Since P(E|S=1) > P(E|S=0), we can infer that P(E=1|S=1) > P(E=1|S=0). However, the fact that P(L|S=1) < P(L|S=0) suggests that the decrease in the probability of being lactose intolerant for smokers is not solely due to the increase in effort.} \\
\\ 
\texttt{In fact, the decrease in the probability of being lactose intolerant for smokers is likely due to the direct effect of smoking on lactose intolerance, which is not specified} \\
            }
        }
    }
    & \\ \\

    \theutterance \stepcounter{utterance}  
    & \multicolumn{4}{p{0.6\linewidth}}{
        \cellcolor[rgb]{0.561,0.737,0.561}{
            \makecell[{{p{\linewidth}}}]{
                \texttt{\tiny{[P1$\rangle$GM]}}
                \texttt{No.} \\
\\ 
\texttt{The direct effect of smoking on lactose intolerance is not specified, but we can infer that it is negative, as the probability of being lactose intolerant for smokers (P(L|S=1)) is less than for nonsmokers (P(L|S=0)). However, the indirect effect of smoking through effort is positive, as effort has a direct effect on lactose intolerance. Since the indirect effect is positive, but the direct effect is not specified, we cannot conclude that smoking decreases the chance of being lactose intolerant.} \\
            }
        }
    }
    & & \\ \\

    \theutterance \stepcounter{utterance}  
    & & & \multicolumn{2}{p{0.3\linewidth}}{
        \cellcolor[rgb]{0.9,0.9,0.9}{
            \makecell[{{p{\linewidth}}}]{
                \texttt{\tiny{[GM$|$GM]}}
                \texttt{No} \\
            }
        }
    }
    & & \\ \\

    \theutterance \stepcounter{utterance}  
    & & & \multicolumn{2}{p{0.3\linewidth}}{
        \cellcolor[rgb]{0.9,0.9,0.9}{
            \makecell[{{p{\linewidth}}}]{
                \texttt{\tiny{[GM$|$GM]}}
                \texttt{yes} \\
            }
        }
    }
    & & \\ \\

    \theutterance \stepcounter{utterance}  
    & & & \multicolumn{2}{p{0.3\linewidth}}{
        \cellcolor[rgb]{0.9,0.9,0.9}{
            \makecell[{{p{\linewidth}}}]{
                \texttt{\tiny{[GM$|$GM]}}
                \texttt{game\_result = LOSE} \\
            }
        }
    }
    & & \\ \\

\end{supertabular}
}

\end{document}
